\documentclass{report}
\usepackage{amsmath,amssymb}
\setlength{\parindent}{0mm}
\setlength{\parskip}{1em}
\begin{document}
\begin{center}
\rule{6in}{1pt} \
{\large Joseph Benzaken \\
{\bf Multigrid Methods for Isogeometric Thin Plate Discretizations}}

2805 Olson Dr Unit B \\ Boulder \\ CO 80303
\\
{\tt Joseph.Benzaken@colorado.edu}\\
Joseph Benzaken\\
Christopher Coley\\
John A. Evans\end{center}

In this talk, we present a novel isogeometric approach to the numerical
solution of the classical Reissner-Mindlin and Kirchhoff-Love plate
equations. This approach eliminates the common issue of locking in thin
plates, a numerical phenomenon resulting from an incompatibility between
the finite element spaces for the translational and rotational
displacement degrees of freedom. This locking-free implementation also
permits the use of a simple geometric multigrid method for solving the
resulting linear system in which Schwarz methods with
intelligently-chosen subdomains are used for iterative smoothing in the
multigrid V-cycles. In the thin plate limit, our multigrid approach
automatically and exactly preserves the constraint that the shear strain
is zero at every geometric level. This results in a method with
convergence rates independent of the thickness. Moreover, this elucidates
the problem as a network of coupled plates with Dirichlet boundary data
specified by adjacent subdomains. Numerical results for both the
Kirchhoff-Love and Reissner-Mindlin plates are presented. The results
demonstrate the robustness of the numerical method through the invariance
of convergence rates with respect to thickness.


\end{document}
